\section{Introduction}

% First graph, Sub-graph pattern counting is important... But privacy issue... To address this, DP has been widely used: by adding... DP guarantee the privacy of node/edge. 



Graph pattern counting is fundamental for extracting meaningful insights from graph data, with applications in community detection~\cite{palla2005uncovering}, fraud detection~\cite{akoglu2013anomaly}, and network characterization~\cite{milo2002network}.
However, directly releasing such analytical results can pose significant privacy risks, as adversaries may infer sensitive information about nodes or edges in the graph~\cite{narayanan2009anonymizing, wu2021adapting, zhang2022inference}.
To address this privacy concern, \emph{differential privacy} (DP)~\cite{dwork2006calibrating} has emerged as a gold-standard framework for privacy-preserving graph data analysis. 
By injecting carefully calibrated noise into the computation, DP prevents the reliable inference of any single node or edge's presence, thereby providing node- or edge-level privacy guarantees.


% However, traditional focuses central... Local model...



Most existing studies on DP in graphs adopt the \emph{central differential privacy} (central-DP) model~\cite{nissim2007smooth,karwa2011private,zhang2015private,dong2023continual}, where a trusted curator accesses the entire dataset and applies the DP mechanism globally.
While effective in utility (i.e., low error), this model assumes full trust in the curator, making it vulnerable to data breaches and privacy attacks. 
To mitigate this risk, the \emph{local differential privacy} (LDP)~\cite{kasiviswanathan2011can,cormode2018privacy} model has gained traction, wherein each data holder privatizes their data locally before sharing it with any external party, including the analyzer and other data holders. 
In the graph setting, each node functions as a data holder, with access only to its incident edges.
Following prior work on LDP-based graph pattern counting~\cite{sun2019analyzing,ye2020towards,ye2020lf,Imola2021,Imola2022,eden2023triangle,he2025robust}, we adopt the same edge-level privacy setting, referred to as edge-LDP.
Consequently, any released information is well protected, offering robust privacy guarantees even under adversarial conditions.


% Following prior work on LDP for graph pattern counting~\cite{}, we focus specifically on protecting edge-level privacy, a setting known as edge-LDP. 



% However, consider... is very complex. Why? a patten requires information of multiple nodes making it complex for privacy protection. As a result.. focuses on specific pattern like star, triangle .... There is no universal solution even only for the general acyclic counting queries. Therefore two challenges: 1. general collect data/ 2.remove node duplication


However, realizing graph pattern counting under LDP is fundamentally challenging due to the decentralized nature of data collection. 
More specifically, constructing a graph pattern typically requires aggregating information distributed across multiple nodes, which is inherently complex under the LDP privacy requirement. 
As a result, existing works have primarily focused on simple, localized structures such as stars~\cite{Imola2021} and triangles~\cite{ye2020towards, ye2020lf,Imola2021, Imola2022,eden2023triangle, he2025robust}, where the required information can be gathered with minimal coordination between nodes.
Despite this progress, no universal solution exists, even for the subclass of acyclic pattern counting queries, which serve as fundamental primitives in network analysis~\cite{milo2002network, itzkovitz2005subgraphs,katz1953new,liben2003link} and relational database queries~\cite{yannakakis1981algorithms, gottlob2001complexity} in the non-private setting.
% LDP motivation
Furthermore, realizing acyclic pattern counting under LDP is natural and practically important when the graph is inherently distributed across multiple parties (e.g., across social platforms).
In this work, we identify the following two unresolved challenges in enabling acyclic graph pattern counting under LDP.




\begin{table}[t]
\centering
% \setlength{\tabcolsep}{5pt}
\renewcommand{\arraystretch}{1.2}
\resizebox{\textwidth}{!}{
\begin{tabular}{l c *{3}{c} *{3}{c}}
\toprule
 &  & \multicolumn{3}{c}{\textbf{Our Results}} & \multicolumn{3}{c}{\textbf{SOTA}} \\
\cmidrule(lr){3-5}\cmidrule(lr){6-8}
\textbf{Task} & \textbf{Query Result}
& \textbf{Error} & \textbf{Comm.} & \textbf{Time}
& \textbf{Error} & \textbf{Comm.} & \textbf{Time} \\

\midrule
$k$-line walk counting
& $\leq N d(G)^{k}$
& $\tilde{O}(\sqrt{N}d(G)^{k-1})$ & $O(M + N)$ & $O(N d(G))$
& 
{-}
% \blue{$O(N^{k/2}d(G)^{k/2})$
% \footnotemark} 
& {$O(N^2)$} & {$O(N d(G))$} \\

$k$-line path counting
& $\leq N d(G)^{k}$
& $\tilde{O}(\sqrt{N} d(G)^{k})$ & $O(M + N)$ & $O(Nd(G))$
& {$O(N^k)$} & $O(N^2)$ & $O(N^{k+1})$ \\


$k$-edge acyclic pattern counting
& $\leq N d(G)^{k}$
& $\tilde{O}(\sqrt{N} d(G)^{k})$ & $O(M + N)$ & $O(Nd(G))$
& {$O(N^k)$} & $O(N^2)$ & $O(N^{k+1})$ \\

% $k$-star counting
% & $\leq N d(G)^{k}$
% & $\tilde{O}(\sqrt{N} d(G)^{k-1})$ & $O(N)$ & $O(N)$
% & $\tilde{O}(\sqrt{N} d(G)^{k-1})$ & $O(N)$ & $O(N)$ \\
\bottomrule
\end{tabular} 
}
\caption{
Summary of results for three pattern counting queries under edge-LDP.\protect\footnotemark\protect\footnotemark
Here, $N$ is the number of nodes, $M$ is the number of edges, $d(G)$ is the maximum degree of graph $G$, and $k$ is the number of edges in the pattern.
SOTA denotes the state-of-the-art method: {\cite{betzer2024publishing} for $k$-line walk counting, and \cite{suppakitpaisarn2025counting} for both $k$-line path counting and $k$-edge acyclic pattern counting.}
{The error bound in \cite{betzer2024publishing} depends on ad hoc parameters and is therefore not included.
}
}
\label{tab:task-comparison}
\end{table}

\footnotetext[1]{
We use $\Omega(\cdot)$ to denote asymptotic lower bounds, $O(\cdot)$ for asymptotic upper bounds, and $\tilde{O}(\cdot)$ to suppress logarithmic factors. 
The number of pattern edges $k$, privacy budget $\varepsilon$, and error probability $\beta$ are treated as constants and omitted from the asymptotic notation.
}

\footnotetext[2]{
All results in the introduction hold with constant probability.
}



\textbf{Challenge 1. Generalizing acyclic pattern construction.} 
% Developing a general method to reconstruct arbitrary acyclic patterns from distributed information.
Existing LDP pattern counting solutions typically exploit structure-specific properties to guide computation, making them hard to extend to more complex or varied patterns.
For example, in star counting, each node independently counts the number of stars centered at itself and sends a perturbed count to the analyzer for aggregation.
In triangle counting, a common approach is for each node to locally enumerate pairs of its neighbors and complete the triangle using the noisy information of the last edge.
In contrast, general acyclic pattern counting should support arbitrary sizes and structures without relying on any pattern-specific assumptions.



\textbf{Challenge 2. Eliminating node duplication.}  
Acyclic patterns require that all nodes within a pattern instance be distinct, since repeated nodes would form cycles.
While this constraint also applies to stars and triangles, duplication is easy to prevent there because all involved nodes lie within a local neighborhood, which allows a single node to ensure that no repeated nodes appear in the pattern. 
In contrast, general acyclic patterns may span distant nodes, making it difficult for any single node to detect and eliminate duplication.





A straightforward approach to addressing both challenges is to process data centrally on the analyzer using the standard LDP technique of \emph{Warner's Randomized Response} (RR)~\cite{warner1965randomized}, {as recently adopted by~\cite{suppakitpaisarn2025counting}}.
Specifically, each node privatizes its incident edge list using RR and sends it to the analyzer, which then reconstructs a global noisy graph and performs analytics on it.
However, this method yields poor utility with an error bound of {$O(N^k)$}, where $N$ is the number of nodes and $k$ is the number of edges in the target pattern. 
In contrast, the true pattern count is at most $N d(G)^k$, where $d(G)$ is the maximum degree. 
{For many real-world networks with relatively uniform degree distributions, the count typically scales as $O(N d(G)^k)$.
This gap highlights the inefficiency of the RR-based method.
In the central-DP literature, algorithms often achieve error bounds of the form $\tilde{O}(\mathrm{poly}(d(G)))$~\cite{johnson2018towards}.
Since LDP typically incurs an additional $O(\sqrt{N})$ overhead~\cite{bonawitz2017practical} compared to central-DP, this raises the central question of our work:}
% For most real-world graphs, prior work~\cite{Imola2021} and our experiments consistently show that the error of the RR approach exceeds the actual count by a wide margin, rendering it impractical.
% This raises the central question of our work:


% \footnotetext[1]{\marginnote{\blue{R2.O2}}\tocomplete{
% The original paper’s error bound depends on ad hoc parameters, and we report the worst-case bound under the ideal clipping threshold $\sqrt{d(G)}$ suggested in that paper.
% }}
% We derive this worst-case error bound under the ideal assumption that the clipping threshold is set to $\sqrt{d(G)}$, as suggested in the original paper.







% !!!!!! DO NOT FORGET:


% As shown in Section~\ref{subsec:strawman}, this method suffers from poor utility with an error of $\Omega(N^{k+1})$
% , and incurs an exponential runtime of $O(N^{k+1})$ on the analyzer side, making it infeasible for large-scale, real-world graphs. 



\emph{{Is there a general solution for counting arbitrary acyclic graph patterns under edge-LDP with an error of $\tilde{O}(\sqrt{N}\,\mathrm{poly}(d(G)))$?}}
% \emph{Is there a practical and general solution for counting arbitrary acyclic graph patterns under edge-LDP?}

% Our result idea is to reduce subproblems.. Each node collect Solve sub-problem and then, based previous ... 
% Give the line path counting query as an example, in each iteration, each collect, from neighbors. To remove ..., we invent rank marking 



This paper answers the question affirmatively. 
To address \textbf{Challenge~1}, we 
%\marginnote{\orange{R3.O2}}\tocomplete{adapt the subtree decomposition~\cite{yannakakis1981algorithms} to the LDP setting}, recursively reducing acyclic pattern counting to subpattern counting and solving it step by step.
recursively reduce acyclic pattern counting to subpattern counting, where each subpattern is also acyclic, and solve it step by step.
Consider path counting as an example, where the goal is to count the number of paths with $k$ edges, namely $k$-line paths. 
In each round, a node collects from its neighbors the counts of paths of a certain length ending at those neighbors, starting from the base case of $0$-line paths, each initialized with a count of $1$. 
The node then aggregates these counts to compute the number of paths that are one edge longer and end at itself.
After $k$ rounds, each node holds the count of $k$-line paths ending at itself, and the global count is obtained by summing over all nodes. 
To address \textbf{Challenge~2}, {we develop a \emph{random marking} technique inspired by color-coding~\cite{alon1995color}, which was originally proposed to reduce the runtime of non-private path and cycle detection.}
 % we propose a novel technique called \emph{random marking}
With random marking, each node is restricted to appear only at a designated position in the pattern during computation. 
This technique can be seamlessly combined with our first solution by allowing each node to participate in only one specific round. 
By integrating both strategies, we develop the first general solution for acyclic pattern counting under edge-LDP. 
% Table~\ref{tab:task-comparison} summarizes the results. 
Our main contributions are as follows:






\begin{enumerate}
\item 
As our first result, we address $k$-line walk counting under edge-LDP, which serves as the foundation of our mechanisms for acyclic patterns. 
Our mechanism runs in $k-1$ rounds and produces an unbiased estimate of the true number of $k$-line walks, with an error of $\tilde{O}(\sqrt{N} d(G)^{k-1})$.
Summing over all rounds, the mechanism incurs an overall communication cost of $O(M+N)$, a computation cost of $O(N)$ on the analyzer, and $O(d(G))$ on each node.


\item
We next address $k$-line path counting by applying random marking to our $k$-line walk counting solution.
The resulting $k$-round mechanism also produces an unbiased estimate, with a slightly increased error of $\tilde{O}(\sqrt{N}d(G)^k)$. 
Moreover, it maintains the same overall communication cost of $O(M+N)$ and computation costs of $O(N)$ on the analyzer and $O(d(G))$ on each node.


\item 
We further generalize our $k$-line path counting solution to handle arbitrary acyclic patterns with $k$ edges. 
The resulting mechanism runs in at most $k$ rounds and produces an unbiased estimate, with asymptotically the same error of $\tilde{O}(\sqrt{N} d(G)^k)$.
The communication and computation costs remain unchanged, with a total communication cost of $O(M+N)$, a computation cost of $O(N)$ on the analyzer, and $O(d(G))$ on each node.
% Additionally, for $k$-star counting, we will show our mechanism can be optimized to achieve an error of $\tilde{O}(\sqrt{N}d(G)^{k-1})$, matching that of the SOTA method.

\end{enumerate}



Experimental results show that our mechanism achieves a relative error below ${8}\%$ on $k$-line walk counting queries. 
For acyclic pattern counting tasks, including $k$-line path counting and $k$-edge acyclic pattern counting, our mechanisms yield up to ${46}$–${2600}\times$ improvement in utility and ${300}$–$650\times$ reduction in communication cost over the RR-based approach. 

% Additionally, for $k$-star counting, our mechanism outperforms $\mathrm{LocalLap}_{k\star}$ with up to a $710\times$ improvement in utility and a $3\times$ reduction in communication cost.





